
%\documentclass[./main.tex(path of main file)]{subfiles}
\documentclass[10pt]{article}
\usepackage{amsmath}
\DeclareMathOperator*{\argmin}{arg\,min} % thin space, limits underneath in displays
\DeclareMathOperator*{\argmax}{arg\,max} % thin space, limits underneath in displays
\newtheorem{thm}{Theorem}
\usepackage{amssymb}
\usepackage{amsfonts}
\usepackage{mathrsfs}
\usepackage{bm}
\usepackage{indentfirst}
\setlength{\parindent}{0em}
\usepackage[margin=1in]{geometry}
\usepackage{graphicx}
\usepackage{setspace}
\doublespacing
\usepackage[flushleft]{threeparttable}
\usepackage{booktabs,caption}
\usepackage{float}
\usepackage[sort,comma]{natbib}
\usepackage[hidelinks]{hyperref}
\usepackage{booktabs}
\usepackage{multirow}

\usepackage{import}
\usepackage{xifthen}
\usepackage{pdfpages}
\usepackage{transparent}
\usepackage{subfiles}
\usepackage{blindtext}
\usepackage{glossaries}
\usepackage{pdflscape}
\usepackage{adjustbox}
\usepackage{rotating}
\usepackage{eurosym,geometry,ulem,subcaption,caption,color,sectsty,comment,footmisc,pdflscape,array}

% Format definition list
\newenvironment{mylist}[1]%
  {\begin{list}{}{\settowidth{\labelwidth}{\textbf{#1}}
   \setlength{\leftmargin}{\labelwidth}
   \addtolength{\leftmargin}{\labelsep}
   \renewcommand{\makelabel}[1]{\textbf{\hfill##1}}}}%
  {\end{list}}



\newcommand{\incfig}[1]{%
\def\svgwidth{\columnwidth}
\import{./figures/}{#1.pdf_tex}
}




\title{}
\author{}
\date{}


\begin{document}
\maketitle

\section{Key concepts}
\subsection{Wealth vs. Income}
Income is a receipt of money by an individual during some period of time, such as labor income, 
interest income.
An individual's wealth is the level of assets (cash, checking accounts, savings accounts, stock, 
bonds, etc.) that an individual has in store. It can be negative if an individual is overdrawn
on his checking account or otherwise is in debt.

{\textbf {Example:}}:\\
If you have \$1000 in your savings accounts, then you have \$1000 in wealth.
If it generates \$30 interest in period 2, then your interest income is \$30.
If you earn a labor income of \$10000 in period 2, your total income is \$10030.

\subsection{First-Order Condition}
A first-order condition with respect to any particular variable mathematically describes how
a maximum is achieved by optimally setting/choosing that particular variable, taking as given the
setting/choices for all other variables.



\section{Dynamic Consumption-Savings Framework}
In this section, we will focus on the study of intertemporal (literally, ``across time")
choices of individuals.
We first ignore  leisure and labor altogether.

Consider a two-period model.

\subsection{Assumptions}
\begin{itemize}
\item 
		Each individual plans economic events for two time periods: present (period 1) and future 
		(period 2). There is no period 3! And every individual knows it!
		We will relax this assumption later.
\item 
		Individual has no control over his or her income. Each individual will receive labor income
		in each of the two periods. 
\item 
		Individual has to make a choice about consumption in each of the two periods.
\item
		Saving or borrowings (dissaving) are allowed during period 1.
\end{itemize}



\subsection{Model}
\subsubsection{Period-by-Period Budget Constraint}

\begin{align}
\text{ Period 1: }P_1c_1 + A_1 &= (1 + i)A_0 + Y_1 \label{eqn:t1_bc}\\
\text{ Period 2: }P_2c_2 + A_2 &= (1 + i)A_1 + Y_2 \label{eqn:t2_bc}
\end{align}
We must have $ A_2 = 0 $ as there is no period 3.



$ Y_1 $ is labor income received at the beginning of period 1. Each individual begins period 1
with some initial wealth (which may be negative), $ A_0 $. 
Each individual chooses consumption $ c_1 $ in period 1, each unit of which costs $ P_1 $ dollars.
He also decides how much wealth to carry into period 2, $ A_1 $. A negative $ A_1 $ means that 
the individual is in debt at the beginning of period 2. Nominal interest rate denotes as $ i $.

Private saving in a given time period is being defined as the difference between an individual's
total income and his total expenditures, consumption and taxes (we ignore taxes for now).
From the above, the total income is the sum of labor income and interest income,
\begin{equation*}
TotalIncome_1 = Y_1 + iA_0.
\end{equation*}
Thus saving in period 1 is
\begin{equation*}
S^{priv} = iA_0 + Y_1 - P_1c_1
\end{equation*}
If we rearrange equation \eqref{eqn:t1_bc},
\begin{equation*}
A_1 - A_0 = iA_0 + Y_1 - P_1c_1,
\end{equation*}
it says saving is the change in wealth during the course of a time period.

\subsubsection{Lifetime Budget Constraint}

$ A_1 $ appears in both equation \eqref{eqn:t1_bc} and \eqref{eqn:t2_bc}. Rearrange equation 
\eqref{eqn:t1_bc} by writting $ A_1 $ as a function of all other terms, then put it into 
equation \eqref{eqn:t2_bc}, we get.
\begin{equation}
P_1c_1 + \frac{P_2c_2}{1 + i} = Y_1 + \frac{Y_2}{(1 + i)} + (1 + i)A_0,
\end{equation}
which is the lifetime budget constraint (LBC).
If we write $ c_2 $ as a function of $ c_1 $, we get
\begin{equation}
c_2 =  - \frac{P_1(1 + i)}{P_2}c_1 + \frac{1 + i}{P_2}Y_2 + \frac{Y_2}{P_2}
\end{equation}







\subsection{Results}
{\textbf {Real Interest Rate and Savings:}}
\begin{itemize}
\item 
		If an individual is initially a debtor (consumption is greater than income) at the end of period 
		1, then a rise in the real interest rate necessarily increases his savings during period 1.
\item 
		If an individual is initially a saver at the end of period 1, then a rise in the real interest
		rate may increase or decrease his savings during period 1.
\end{itemize}
{\textbf {Yet theory cannot guide use as to how private savings at the macroeconomic level
responds to a rise in the real interest rate!}}
Many empirical studies conclude that the real interest rate in fact has a {\underline {very weak 
effect}}, if any effect at all, on private savings behavior.

The studies that do show that real interest rates do influence savings almost always conclude that
a rise in the real interest rate leads to a rise in savings.



\bibliographystyle{plainnat}
\bibliography{my_bib}

\end{document}

